\chapter{The Translation Layer}

\epigraph{"Clients often need to have a conversation with someone who thinks like a mathematician: carefully defining terms; making implicit assumptions explicit, and knowing when an approximation is or is not good enough."}{SIAM News interview (2022)}

Influence isn't in pageviews. It's in production code.

Evidence: citations in source files, implementations ported across languages, official documentation linking to a WordPress blog.

\bigskip

\section*{The pandas Citation}

In the source code of pandas, the Python library used by millions of data scientists, with 47,000+ GitHub stars, there's a file called \texttt{groupby.pyx}. It contains this comment:

\begin{lstlisting}[language=python]
# Running stats update based on RunningStats::Push from
# https://www.johndcook.com/blog/skewness_kurtosis/
\end{lstlisting}

Every time a data scientist calls \texttt{.skew()} or \texttt{.kurt()} on a grouped DataFrame, this code runs. Cook's running statistics algorithm, based on a 1962 paper by B. P. Welford, computes these higher moments in a numerically stable way.

The naive formula for variance (E[X$^2$] - E[X]$^2$) fails catastrophically with large numbers. Welford's algorithm doesn't. Cook explained it, implemented it. Pandas developers chose it.

Infrastructure. Not famous. Not attributed in user-facing documentation. Running silently in data pipelines worldwide.

\bigskip

\section*{The OpenTSDB Citation}

OpenTSDB is a time-series database used for monitoring and metrics at scale (5,000+ GitHub stars). In its core file \texttt{Aggregators.java}:

\begin{lstlisting}[language=java]
/**
\begin{itemize}
\item This implementation is based upon a paper by John D. Cook,
\item which itself is based upon a method that goes back to a 1962
\item paper by B. P. Welford
\end{itemize}
 */
\end{lstlisting}

Same algorithm. Different language. Same attribution.

When a company monitors its servers using OpenTSDB, every variance calculation traces back to Cook's blog.

\bigskip

\section*{The TensorFlow Documentation}

The official TensorFlow documentation includes a guide on automatic differentiation. In that guide:

\begin{blogquote}
Note: Don't actually invert the matrix.
\end{blogquote}

This is Google's machine learning framework, used by millions of developers, linking to a WordPress blog post for a key concept.

The post explains why matrix inversion is computationally wasteful. If you need to solve $Ax = b$, you should use LU decomposition or Cholesky factorization, not compute $A^{-1}$. Textbooks often write $x = A^{-1}b$, which is mathematically correct but computationally disastrous.

Cook's post became the canonical reference. TensorFlow's documentation doesn't explain the concept itself---it links to Cook.

\bigskip

\section*{The Julia Documentation}

The official Julia language documentation for floating-point numbers includes:

\begin{blogquote}
For a brief but lucid presentation of how floating-point numbers are represented, see John D. Cook's article...
\end{blogquote}

Julia is a language designed for scientific computing (48,000+ GitHub stars). Its documentation recommends Cook for understanding IEEE 754.

\bigskip

\section*{The Ten-Language Port}

Cook's gamma function implementation has been ported to at least ten programming languages:

\begin{longtable}{l p{6.5cm}}
\toprule
\textbf{Language} & \textbf{Repository} \tabularnewline
\midrule
C\# & ghewgill/picomath \tabularnewline
Java & raharrison/MathEngine \tabularnewline
JavaScript & AndreasMadsen/mathfn \tabularnewline
C++ & Vaa3D/vaa3d\_tools \tabularnewline
PHP & MatiasNAmendola/php-nlp-tools \tabularnewline
PowerShell & gravejester/Communary.ToolBox \tabularnewline
Scheme & shirok/Gauche \tabularnewline
TypeScript & roualdes/distributions \tabularnewline
Rust & Eh2406/quadrature \tabularnewline
D & p0nce/turtle \tabularnewline
\bottomrule
\end{longtable}

Each port cites johndcook.com. The gamma function---which extends the factorial to non-integers---is needed across all mathematical computing. Cook's implementation, using the Lanczos approximation, became the template.

Greg Hewgill created picomath specifically to port Cook's code:

\begin{blogquote}
"John D. Cook is the original author of most of the functions themselves."
\end{blogquote}

\bigskip

\section*{Python's Standard Library}

Python Enhancement Proposal 450 (PEP 450), which added the \texttt{statistics} module to Python's standard library, cites Cook's work on computing standard deviation.

Not a third-party library. Python itself. The \texttt{statistics} module that ships with every Python installation was informed by Cook's blog.

\bigskip

\section*{Five Examples}

How explanations become infrastructure:

Textbooks often present opaque formulations. The naive variance formula, the integral definition of the Gamma function:

\begin{itemize}
\item   \textbf{Welford Naive Variance:}
\end{itemize}
    `$$s\textasciicircum{}2 = \frac{\sum (x\_i - \bar{x})\textasciicircum{}2}{N-1}$$`
\begin{itemize}
\item   \textbf{Gamma Function Integral:}
\end{itemize}
    `$$\\Gamma(z) = \int\_0\textasciicircum{}\\infty x\textasciicircum{}{z-1} e\textasciicircum{}{-x} dx$$`

Cook translates these into clarity and working code.

\textbf{1. "Don't Invert That Matrix" (2010)}

Textbooks write $x = A^{-1}b$. Students read it literally. Nobody explains computational cost.

Cook names the problem, provides alternatives (LU, Cholesky), explains why in 500 words.

Now in TensorFlow documentation. QuantEcon lectures. 25+ GitHub repositories.

\textbf{2. Gamma Function Approximation}

$\Gamma(z)$ defined as an improper integral---useless for computation. Textbooks don't explain implementation.

Cook provides Lanczos approximation code, explains when to use log-gamma, handles edge cases.

Now in pandas core. Ten language implementations. Microsoft Azure libraries.

\textbf{3. Running Statistics (Welford's Algorithm)}

Naive variance formula fails numerically. Welford's 1962 paper behind a paywall and dense.

Cook explains the algorithm, provides code, shows why naive formula fails.

Now in pandas groupby. OpenTSDB. Multiple time-series libraries.

\textbf{4. Kumaraswamy Distribution (2017)}

Beta distributions expensive to sample---no closed-form inverse CDF. Kumaraswamy's 1980 paper exists but doesn't reach practitioners.

Cook explains the comparison, provides closed-form expressions and code.

Hacker News: "I was looking for ways to sample quickly from a mixture of Beta distributions when I stumbled upon his post... makes generating samples convenient and fast."

\textbf{5. RGB to Grayscale (2009)}

Libraries provide conversion but don't explain which algorithm or why coefficients exist.

Cook explains perceptual luminance, vision science behind "magic numbers," when each method is appropriate.

First Google result for "RGB to grayscale." HN: "I just googled 'RGB to grayscale' and this was the first result."

\bigskip

\section*{The Pattern}

Academic source exists but inaccessible (behind paywall, too dense, wrong format). Cook translates into clear prose plus working code. Developers adopt because the explanation is findable and usable. Infrastructure forms as implementations spread.

Cook didn't invent Welford's algorithm or the Lanczos approximation. He translated them from academic papers into code that works. Translation is contribution.

\bigskip

\section*{The Quantified Impact}

\begin{longtable}{l p{6.5cm}}
\toprule
\textbf{Metric} & \textbf{Value} \tabularnewline
\midrule
GitHub repos citing johndcook.com & 160+ \tabularnewline
Combined stars of citing repos & 500,000+ \tabularnewline
Languages with gamma implementation & 10+ \tabularnewline
Official documentation citations & 3 (TensorFlow, Julia, PEP 450) \tabularnewline
\bottomrule
\end{longtable}

These numbers undercount. Only explicit citations. Many codebases adapt without attribution. Many developers learn and apply without linking back.

\bigskip

\section*{The Translation Layer Thesis}

Cook occupies a rare position: PhD-level mathematical knowledge, decades of programming experience, and the ability to explain clearly.

\begin{longtable}{l p{6.5cm}}
\toprule
\textbf{Credential} & \textbf{Evidence} \tabularnewline
\midrule
Academic depth & PhD mathematics, 5,881 Google Scholar citations \tabularnewline
Industry experience & 30+ years; clients include Google, Amazon, Microsoft, Facebook, AMD \tabularnewline
Communication skill & 5,233 posts, 500,000+ Twitter followers \tabularnewline
\bottomrule
\end{longtable}

Most mathematicians can't program fluently. Most programmers can't derive the gamma function. Most writers can't do either. Cook can do all three.

This intersection is the translation layer: academic mathematics into production code. Textbook definitions into working implementations.

\bigskip

\section*{What Readers Say}

\textbf{abhgh on Hacker News:}
\begin{blogquote}
"I have personally benefited from it; I was looking for ways to sample quickly from a mixture of Beta distributions when I stumbled upon his post on the Kumaraswamy distribution."
\end{blogquote}

\textbf{Tim Hopper (data scientist):}
\begin{blogquote}
"John's writings have shaped many aspects of my trajectory from being an operations research grad student to a practicing data scientist."
\end{blogquote}

\textbf{Anna Haensch (AMS Blog on Math Blogs):}
\begin{blogquote}
"I always walk away feeling like I learned something."
\end{blogquote}

\textbf{xpil:}
\begin{blogquote}
"Your blog is one of the most valuable pieces of the Internet."
\end{blogquote}

\bigskip

\section*{Infrastructure, Not Fame}

Nearly 6,000 Google Scholar citations. Over 1,200 Hacker News submissions. University syllabi at MIT, Stanford, and Carnegie Mellon.

But not famous. No keynotes. No bestselling book. Not verified on Twitter.

Infrastructural influence. Living in code comments, production systems, the understanding of thousands of developers who needed exactly that explanation at 2 AM.

Quiet influence. Not followers. Not fame. Usefulness, compounded over seventeen years.

\bigskip

\textit{The evidence is in production code. The impact is infrastructural. But what about everything the data can't capture?}

